\documentclass[10pt]{article}
\usepackage[a4paper,margin=0.86in]{geometry}
\renewcommand{\rmdefault}{cmr}
\usepackage{amsmath,amssymb,amsthm,mathtools,amscd,enumitem,xcolor,hyperref}
\hypersetup{colorlinks=true,linkcolor=blue!55!black,citecolor=blue!55!black,urlcolor=blue!55!black}
\setlength{\emergencystretch}{4em}
\setlength{\parindent}{0pt}
\setlength{\parskip}{0.52em}
\setlist{leftmargin=1.55em,itemsep=0.2em,topsep=0.25em}
\sloppy

\newtheorem{theorem}{Theorem}[section]
\newtheorem{proposition}[theorem]{Proposition}
\newtheorem{lemma}[theorem]{Lemma}
\newtheorem{corollary}[theorem]{Corollary}
\theoremstyle{definition}
\newtheorem{definition}[theorem]{Definition}
\newtheorem{convention}[theorem]{Convention}
\newtheorem{remark}[theorem]{Remark}

\newcommand{\Zl}{\mathbf Z_\ell}
\newcommand{\Ql}{\mathbf Q_\ell}
\newcommand{\Qlb}{\overline{\mathbf Q}_\ell}
\newcommand{\OO}{\mathcal O}
\newcommand{\Dbc}{D^b_c}
\newcommand{\Dctf}{D^b_{\mathrm{ctf}}}
\newcommand{\RHom}{R\mathcal Hom}
\newcommand{\RG}{R\Gamma}
\newcommand{\RGc}{R\Gamma_c}
\newcommand{\DD}{\mathbb D}
\newcommand{\Tr}{\operatorname{Tr}}
\newcommand{\red}{\operatorname{red}}
\newcommand{\DctfW}{D^b_{\mathrm{ctf,W}}}

\begin{document}
\begin{center}
{\large\bfseries The Poincare Pairing\par}
\medskip
{\normalsize Compact support, perfectness, degree indexing, and Frobenius reciprocity\par}
\end{center}

\begin{abstract}
For a smooth separated \(d\)-dimensional scheme, Verdier duality and smooth
purity produce a perfect pairing between compactly supported cohomology and
ordinary cohomology.  This note gives the complete chain of definitions and
intermediate statements: global sections, compact support, point
perfectness, cup product, evaluation, trace, derived perfectness,
cohomological indexing, Frobenius similitude, and the proper and relative
smooth-proper forms.
\end{abstract}

\tableofcontents

\section{Coefficient categories and perfect complexes at a point}

Fix a prime \(\ell\), a separably closed field \(k\) of characteristic
different from \(\ell\), and a separated finite-type \(k\)-scheme \(X\).

\begin{definition}[Coefficient stages]
The coefficient ring \(A\) is one of the following:

\begin{enumerate}
\item a finite local \(\mathbf Z/\ell^n\mathbf Z\)-algebra \(\Lambda\);
\item the ring of integers \(\OO_E\) of a finite extension \(E/\Ql\);
\item the field \(E\);
\item \(\Qlb\), interpreted as the Karoubi-completed filtered \(2\)-colimit
over finite \(E/\Ql\).
\end{enumerate}

For every finite extension \(E/\Ql\), choose a uniformizer \(\varpi\) of
\(\OO_E\) and put

\[
 \Lambda_n=\OO_E/\varpi^n\qquad(n\geq1).
\]

For finite \(A=\Lambda\), write

\[
 D(X,A)=D\bigl(\mathsf{Mod}(X_{\mathrm{pro\acute et}},A_X)\bigr).
\]

The category \(\Dbc(X,\Lambda)\) has bounded constructible cohomology, and
\(\Dctf(X,\Lambda)\) is its finite-Tor subcategory.  Define
\(D(X,\OO_E)\) to be the ordinary triangulated category of normalized systems
\(K=(K_n)_{n\geq1}\), where \(K_n\in D(X,\Lambda_n)\), together with coherent
isomorphisms
\[
 K_{n+1}\otimes^L_{\Lambda_{n+1,X}}\Lambda_{n,X}
 \xrightarrow{\sim}K_n.
\]
The full subcategory \(\Dbc(X,\OO_E)\) is formed by systems having common
cohomological bounds and a common finite constructible stratification.  Its
full subcategory \(\Dctf(X,\OO_E)\) consists of systems with one common finite
Tor-amplitude interval.  Put
\[
 D(X,E)=\operatorname{Kar}\bigl(D(X,\OO_E)[\varpi^{-1}]\bigr),
\]
where \([\varpi^{-1}]\) is the ordinary triangulated localization in which
multiplication by \(\varpi\) is invertible and \(\operatorname{Kar}\) denotes
idempotent completion.  Define \(\Dbc(X,E)\) and \(\Dctf(X,E)\) as the full subcategories
represented by objects of \(\Dbc(X,\OO_E)\) and \(\Dctf(X,\OO_E)\),
respectively.  Finally set
\[
 D(X,\Qlb)=
 \operatorname{Kar}\!\left(
   2\!\operatorname*{colim}_{E/\Ql\ \mathrm{finite}}D(X,E)
 \right),
\]
with transition functors given by finite extension of coefficient fields.
The categories \(\Dbc(X,\Qlb)\) and \(\Dctf(X,\Qlb)\) are the full
subcategories represented by bounded constructible and finite-Tor objects
over some finite coefficient field.  Every category in this paragraph is an
ordinary triangulated category.
\end{definition}

\begin{definition}[Finite Tor-amplitude]
For finite \(A=\Lambda\) and \(K\in D(X,\Lambda)\), say that \(K\) has
Tor-amplitude in \([a,b]\) if for
every sheaf \(M\) of \(A_X\)-modules concentrated in degree zero,

\[
 H^r(K\otimes^L_{A_X}M)=0\qquad(r\notin[a,b]).
\]

The category \(\Dctf(X,\Lambda)\) is the full subcategory of
\(\Dbc(X,\Lambda)\) formed by objects having some finite Tor-amplitude
interval.  At the adic stage one requires one interval valid at every finite
level.  At the \(E\)- and \(\Qlb\)-stages, finite Tor-amplitude means that a
representative over an integral lattice has such a common interval; this
condition is independent of further finite coefficient extension and of
idempotent splitting.
\end{definition}

\begin{definition}[Derived internal Hom and perfect objects on a scheme]
At a finite coefficient stage \(A=\Lambda\), choose a functorial K-flat
replacement \(P(K)\to K\) and a functorial K-injective replacement
\(L\to I(L)\).  Define

\[
 \RHom_X(K,L)=
 \mathcal Hom_X^\bullet(P(K),I(L))\in D(X,A),
\]

where

\[
 \mathcal Hom_X^n(P,I)=
 \prod_r\mathcal Hom_{A_X}(P^r,I^{r+n}),
 \qquad d(\phi)=d_I\phi-(-1)^n\phi d_P.
\]

An object of \(D(X,\Lambda)\) is perfect over \(\Lambda_X\) if it is locally
quasi-isomorphic to a bounded complex of finite projective
\(\Lambda_X\)-modules.  Perfect objects have finite Tor-amplitude, are dualizable,
and the derived tensor--Hom adjunction is

\[
 \operatorname{Hom}(P\otimes^LK,L)
 \cong\operatorname{Hom}(P,\RHom_X(K,L)).
\]

For normalized adic systems with finite-Tor first variable, derived internal
Hom is the normalized system of the finite-level objects
\(\RHom_X(K_n,L_n)\), using the reduction--internal-Hom exchange
isomorphisms.  Localization at \(\varpi\), finite extension of coefficients,
and restriction to an idempotent summand define \(\RHom_X\) at the \(E\)- and
\(\Qlb\)-stages.  Perfectness at these stages means that the object has a
perfect representative at every finite level with one common amplitude.
\end{definition}

\begin{definition}[Finite reduction]
For \(m\leq n\), set

\[
 \red_{n,m,X}(K)=K\otimes^L_{\Lambda_{n,X}}\Lambda_{m,X},
 \qquad
 \red_{n,m,X}:\Dctf(X,\Lambda_n)\longrightarrow
 \Dctf(X,\Lambda_m).
\]

The canonical identifications
\(\red_{m,r}\red_{n,m}\cong\red_{n,r}\) satisfy the cocycle equation.
For a normalized \(\OO_E\)-system \(K=(K_n)\), write
\(\red_nK=K_n\).
\end{definition}

\begin{definition}[Derived category at the point]
Let \(s=\operatorname{Spec}k\).  The pro-etale derived category
\(D(s,A)\), on bounded constructible objects, is identified with the
ordinary derived category \(D(A)\) of \(A\)-modules.  For \(A=\Qlb\), this
means the corresponding Karoubi-completed filtered \(2\)-colimit.
\end{definition}

\begin{definition}[Perfect point complex]
An object \(C\in D(A)\) is perfect if it is quasi-isomorphic to a bounded
complex of finite projective \(A\)-modules.  Its derived dual is

\[
 C^\vee=R\operatorname{Hom}_A(C,A).
\]

If \(A\) is a field, perfectness is equivalent to boundedness and
finite-dimensionality of all cohomology groups.
\end{definition}

\begin{proposition}[Biduality at the point]
For a perfect point complex \(C\), evaluation and coevaluation give

\[
 C\xrightarrow{\sim}R\operatorname{Hom}_A(C^\vee,A),
\]

and for every \(D\in D(A)\),

\[
 C^\vee\otimes_A^LD\xrightarrow{\sim}R\operatorname{Hom}_A(C,D).
\]
\end{proposition}

\section{Ordinary and compactly supported global sections}

Let \(a_X:X\to s\) be the structural morphism.

\begin{definition}[Finite-stage ordinary derived global sections]
Let \(A=\Lambda\) be finite.  Choose a functorial K-injective replacement
\(K\to I(K)\).  Define

\[
 \RG(X,K)=Ra_{X*}K=a_{X*}I(K)\in D(A),
\]

and

\[
 H^i(X,K)=H^i(\RG(X,K)).
\]
\end{definition}

\begin{definition}[Finite-stage compactly supported global sections]
Let \(A=\Lambda\) be finite.  For a separated finite-type structural
morphism, choose a compactification

\[
 X\xrightarrow j\overline X\xrightarrow{\bar a}s
\]

with \(j\) open and \(\bar a\) proper.  Define

\[
 \RGc(X,K)=Ra_{X!}K=R\bar a_*j_!K\in D(A),
\]

and

\[
 H_c^i(X,K)=H^i(\RGc(X,K)).
\]
\end{definition}

\begin{theorem}[Independence and functoriality of compact support]
For finite coefficients, the following assertions hold.
Common refinements of compactifications give canonical isomorphisms between
the corresponding objects \(R\bar a_*j_!K\).  They satisfy identity and
cocycle equations and define an exact triangulated functor

\[
 \RGc(X,-):D(X,A)\longrightarrow D(A).
\]

If \(X\) is proper, the natural comparison

\[
 \RGc(X,K)\xrightarrow{\sim}\RG(X,K)
\]

is an isomorphism.
\end{theorem}

\begin{theorem}[Bounded constructibility of global sections]
For finite \(A=\Lambda\) and \(K\in\Dctf(X,\Lambda)\), both

\[
 \RG(X,K),\qquad\RGc(X,K)
\]

belong to the bounded constructible category at the point.  Their
cohomological amplitudes are bounded in terms of the amplitude of \(K\) and
\(\dim X\).
\end{theorem}

\begin{theorem}[Point perfectness at finite coefficients]
For finite \(A=\Lambda\) and \(K\in\Dctf(X,\Lambda)\), the complexes

\[
 \RG(X,K),\qquad\RGc(X,K)
\]

are perfect objects of \(D(\Lambda)\).
\end{theorem}

\begin{theorem}[Adic and field global-sections functors]
For \(K=(K_n)\in\Dctf(X,\OO_E)\), finite reduction gives coherent
isomorphisms
\[
 \red_{n+1,n}\RG(X,K_{n+1})\xrightarrow{\sim}\RG(X,K_n),
 \qquad
 \red_{n+1,n}\RGc(X,K_{n+1})\xrightarrow{\sim}\RGc(X,K_n).
\]
The resulting normalized systems define exact triangulated functors
\[
 \RG(X,-),\ \RGc(X,-):\Dctf(X,\OO_E)\longrightarrow D(\OO_E).
\]
Localizing at \(\varpi\) defines the corresponding functors to \(D(E)\), and
finite coefficient extension followed by the filtered \(2\)-colimit and
idempotent splitting defines the functors to \(D(\Qlb)\).  These constructions
are the definitions of \(\RG\) and \(\RGc\) at the nonfinite coefficient
stages.  They preserve boundedness, constructibility at the point, and finite
Tor-amplitude.
\end{theorem}

\begin{theorem}[Point perfectness at adic and field stages]
Let \(K\in\Dctf(X,\OO_E)\).  Every reduction of
\(\RG(X,K)\) and \(\RGc(X,K)\) is perfect over
\(\Lambda_n=\OO_E/\varpi^n\), with one common amplitude.  Derived
completeness gives perfect \(\OO_E\)-complexes.  Rationalization gives
perfect \(E\)-complexes, and idempotent completion gives perfect
\(\Qlb\)-complexes.
\end{theorem}

\section{Tensor products, cup products, and evaluation}

\begin{definition}[Coefficient dual on \(X\)]
For \(K\in\Dctf(X,A)\), define

\[
 K^\vee=\RHom_X(K,A_X).
\]

If \(K\) is perfect over \(A_X\), evaluation is

\[
 \operatorname{ev}_K:
 K\otimes^L_{A_X}K^\vee\longrightarrow A_X,
\]

and \(K\to(K^\vee)^\vee\) is an isomorphism.
\end{definition}

\begin{definition}[Mixed-support cup product]
For \(K,L\in D(X,A)\), the projection formula and the counit
\(La_X^*Ra_{X*}L\to L\) define

\[
 \smile_c:
 \RGc(X,K)\otimes_A^L\RG(X,L)
 \longrightarrow
 \RGc(X,K\otimes^L_{A_X}L).
\]

The compactly supported factor is the first factor.
\end{definition}

\begin{proposition}[Cup-product coherence]
The mixed-support cup product is natural in both variables, associative,
unital, and compatible with the symmetry constraint

\[
 x\otimes y\longmapsto(-1)^{|x||y|}y\otimes x.
\]

It is compatible with distinguished triangles, compactification comparison,
finite reduction, and scalar extension.  If \(X\) is proper, it
is the ordinary derived cup product.
\end{proposition}

\section{Dualizing complex and trace}

\begin{definition}[Dualizing complex]
The right adjoint of \(Ra_{X!}\) is

\[
 Ra_X^!:D(A)\longrightarrow D(X,A).
\]

Define

\[
 \omega_X=Ra_X^!A.
\]

The Verdier dual of \(K\in\Dctf(X,A)\) is

\[
 \DD_XK=\RHom_X(K,\omega_X).
\]
\end{definition}

\begin{definition}[Exceptional trace]
The counit of \(Ra_{X!}\dashv Ra_X^!\) is

\[
 \Tr_X:\RGc(X,\omega_X)=Ra_{X!}Ra_X^!A\longrightarrow A.
\]
\end{definition}

\begin{definition}[Verdier evaluation]
For \(K\in\Dctf(X,A)\), tensor--Hom adjunction supplies

\[
 \operatorname{ev}^{\mathrm V}_K:
 K\otimes^L_{A_X}\DD_XK
 =K\otimes^L\RHom_X(K,\omega_X)\longrightarrow\omega_X.
\]
\end{definition}

\begin{definition}[Global Verdier pairing]
Define \(b^{\mathrm V}_K\) to be the typed composite

\[
\begin{aligned}
 \RGc(X,K)\otimes_A^L\RG(X,\DD_XK)
 &\xrightarrow{\smile_c}
 \RGc(X,K\otimes^L\DD_XK)\\
 &\xrightarrow{\RGc(\operatorname{ev}^{\mathrm V}_K)}
 \RGc(X,\omega_X)
 \xrightarrow{\Tr_X}A.
\end{aligned}
\]
\end{definition}

\begin{definition}[Adjoint global-duality comparison]
The tensor--Hom adjoint of \(b^{\mathrm V}_K\) is

\[
 \theta_{X,K}:
 \RG(X,\DD_XK)
 \longrightarrow
 R\operatorname{Hom}_A(\RGc(X,K),A).
\]
\end{definition}

\begin{theorem}[Exceptional tensor--Hom comparison at the point]
For \(K\in\Dctf(X,A)\), the projection formula, the adjunction
\(Ra_{X!}\dashv Ra_X^!\), and internal tensor--Hom adjunction define a
canonical isomorphism

\[
 Ra_{X*}\RHom_X(K,Ra_X^!A)
 \xrightarrow{\sim}
 R\operatorname{Hom}_A(Ra_{X!}K,A).
\]

Its source is \(\RG(X,\DD_XK)\) and its target is
\(R\operatorname{Hom}_A(\RGc(X,K),A)\).
\end{theorem}

\begin{proposition}[Identification of the comparison]
Under the preceding identifications, the exceptional tensor--Hom
isomorphism is exactly \(\theta_{X,K}\), the adjoint of the global Verdier
pairing.  This identification is natural in \(K\), compatible with
distinguished triangles, and compatible with proper compactification.
\end{proposition}

\begin{corollary}[Perfectness of the global Verdier pairing]
For \(K\in\Dctf(X,A)\), the complexes \(\RGc(X,K)\) and
\(\RG(X,\DD_XK)\) are perfect point complexes, and

\[
 b^{\mathrm V}_K:
 \RGc(X,K)\otimes_A^L\RG(X,\DD_XK)\longrightarrow A
\]

is a perfect derived pairing.  Its tensor--Hom adjoint is the isomorphism
\(\theta_{X,K}\).
\end{corollary}

\section{Smooth purity and the explicit Poincare pairing}

Assume from now on that \(X\) is smooth of pure dimension \(d\) over \(k\).

\begin{definition}[Tate twists]
For finite coefficients,

\[
 A_X(1)=\mu_{\ell^n,X}
\]

after the evident scalar extension.  The adic Tate line is the normalized
system with transition \(\zeta\mapsto\zeta^\ell\); rational and \(\Qlb\)
Tate lines are obtained by scalar extension and idempotent completion.  Put

\[
 K(r)=K\otimes^L_{A_X}A_X(r)
\]

for the tensor power of the Tate line when \(r\geq0\) and its tensor inverse
when \(r<0\).
\end{definition}

\begin{theorem}[Smooth purity for the structural morphism]
There is a canonical isomorphism

\[
 \omega_X=Ra_X^!A\xrightarrow{\sim}A_X(d)[2d].
\]

It is compatible with etale localization, products, extension of the
separably closed base field, finite reduction, and scalar extension.
\end{theorem}

\begin{definition}[Normalized smooth trace]
Using smooth purity, define

\[
 \Tr_X^{\mathrm{sm}}:
 \RGc(X,A_X(d)[2d])\xrightarrow{\sim}
 \RGc(X,\omega_X)\xrightarrow{\Tr_X}A.
\]

On cohomology it induces

\[
 \operatorname{tr}_X:
 H_c^{2d}(X,A_X(d))\longrightarrow A.
\]

The normalization sends the ordered product of the standard affine-line
orientation classes on \(\mathbf A^d\) to \(1\).
\end{definition}

\begin{proposition}[Verdier dual of a perfect complex on a smooth scheme]
If \(K\) is perfect over \(A_X\), smooth purity and tensor duality give

\[
 \DD_XK
 =\RHom_X(K,\omega_X)
 \xrightarrow{\sim}K^\vee(d)[2d].
\]
\end{proposition}

\begin{definition}[Derived Poincare pairing]
For a perfect \(K\in\Dctf(X,A)\), define

\[
 \beta_K:
 \RGc(X,K)\otimes_A^L\RG(X,K^\vee(d))[2d]\longrightarrow A
\]

as the composite

\[
\begin{aligned}
 \RGc(X,K)\otimes_A^L\RG(X,K^\vee(d))[2d]
 &\xrightarrow{\smile_c}
 \RGc(X,K\otimes^LK^\vee(d))[2d]\\
 &\xrightarrow{\operatorname{ev}_K}
 \RGc(X,A_X(d))[2d]
 \xrightarrow{\Tr_X^{\mathrm{sm}}}A.
\end{aligned}
\]
\end{definition}

\begin{definition}[Perfect derived pairing]
Let \(C,D\in D(A)\) be perfect.  A morphism

\[
 b:C\otimes_A^LD\longrightarrow A
\]

is perfect if its adjoint

\[
 D\longrightarrow R\operatorname{Hom}_A(C,A)
\]

is an isomorphism.  Equivalently, the adjoint in the first variable is an
isomorphism.
\end{definition}

\begin{theorem}[Derived Poincare duality]
Let \(X/k\) be smooth, separated, and of pure dimension \(d\), and let
\(K\in\Dctf(X,A)\) be perfect over \(A_X\).  Then \(\beta_K\) is a perfect
derived pairing.  Its adjoint is

\[
 \RG(X,K^\vee(d))[2d]
 \xrightarrow{\sim}
 R\operatorname{Hom}_A(\RGc(X,K),A).
\]
\end{theorem}

\section{Cohomological indexing}

\begin{proposition}[Degree of the pairing]
The shift convention \(H^r(C[m])=H^{r+m}(C)\) identifies the degree-zero
part of \(\beta_K\) with pairings

\[
 \langle-,-\rangle_i:
 H_c^i(X,K)\times
 H^{2d-i}(X,K^\vee(d))\longrightarrow A.
\]

Before inserting the Tate twist, the same morphism is the line-valued
pairing

\[
 H_c^i(X,K)\times H^{2d-i}(X,K^\vee)
 \longrightarrow A(-d).
\]
\end{proposition}

\begin{theorem}[Degreewise perfectness over a field]
If \(A\) is \(E\) or \(\Qlb\), all cohomology spaces in the preceding display
are finite-dimensional and each map

\[
 H^{2d-i}(X,K^\vee(d))
 \xrightarrow{\sim}H_c^i(X,K)^\vee
\]

is an isomorphism.  Equivalently, each bilinear pairing
\(\langle-,-\rangle_i\) is perfect.
\end{theorem}

\begin{proposition}[Integral degreewise form]
For \(A=\OO_E\), derived perfectness yields the universal-coefficient exact
sequences for the cohomology modules.  If the relevant cohomology modules
are finite projective over \(\OO_E\), the Ext terms vanish and the
degreewise maps are perfect \(\OO_E\)-module pairings.
\end{proposition}

\begin{corollary}[Constant coefficients]
For \(K=A_X\), Poincare duality is

\[
 H_c^i(X,A)\times H^{2d-i}(X,A(d))\longrightarrow A.
\]

If \(X\) is geometrically connected, then

\[
 H_c^{2d}(X,A(d))\cong A,
 \qquad
 H_c^{2d}(X,A)\cong A(-d),
\]

with the first isomorphism induced by the trace.
\end{corollary}

\section{Proper and smooth-proper specializations}

\begin{definition}[Finite-stage derived direct image]
Let \(A=\Lambda\) be finite.  For a morphism \(f:X\to Y\), choose a
functorial K-injective replacement \(K\to I(K)\) and set

\[
 Rf_*K=f_*I(K),\qquad
 Rf_*:D(X,A)\longrightarrow D(Y,A).
\]

The adjunction \(Lf^*\dashv Rf_*\) has counit
\(Lf^*Rf_*K\to K\).
\end{definition}

\begin{theorem}[Proper constructibility, finite Tor-amplitude, and base change]
Let \(A=\Lambda\) be finite and let \(f:X\to Y\) be proper and of finite
presentation.  Then

\[
 Rf_*:\Dbc(X,A)\longrightarrow\Dbc(Y,A),
 \qquad
 Rf_*:\Dctf(X,A)\longrightarrow\Dctf(Y,A).
\]

For every Cartesian square

\[
\begin{CD}
X' @>{u}>> X\\
@V{f'}VV @VV{f}V\\
Y' @>{v}>> Y,
\end{CD}
\]

the proper base-change map

\[
 Lv^*Rf_*K\xrightarrow{\sim}Rf'_*Lu^*K
\]

is an isomorphism for \(K\in\Dbc(X,A)\), and it preserves the finite-Tor
subcategory.  These maps satisfy identity, composition, and pasting
coherence.
\end{theorem}

\begin{theorem}[Coefficient transport of proper direct image]
Let \(f:X\to Y\) be proper and of finite presentation.  For every normalized
system \(K=(K_n)\in\Dctf(X,\OO_E)\), the finite-stage exchange maps
\[
 \red_{n+1,n,Y}Rf_*K_{n+1}\xrightarrow{\sim}
 Rf_*\red_{n+1,n,X}K_{n+1}\xrightarrow{\sim}Rf_*K_n
\]
make \((Rf_*K_n)_n\) a normalized system.  They define exact functors
\[
 Rf_*:\Dbc(X,\OO_E)\longrightarrow\Dbc(Y,\OO_E),
 \qquad
 Rf_*:\Dctf(X,\OO_E)\longrightarrow\Dctf(Y,\OO_E).
\]
Localization at \(\varpi\), finite extension of coefficient fields, the
filtered \(2\)-colimit, and idempotent splitting define the corresponding
functors for \(A=E\) and \(A=\Qlb\).  At every coefficient stage their source
and target are the displayed bounded constructible and finite-Tor categories,
and the proper base-change maps are isomorphisms compatible with coefficient
transport.
\end{theorem}

\begin{definition}[Relative cup product]
For \(K,L\in D(X,A)\), the lax monoidal structure on \(Rf_*\), equivalently
the two adjunction counits after pulling back to \(X\), defines

\[
 \smile_f:Rf_*K\otimes^L_{A_Y}Rf_*L
 \longrightarrow Rf_*(K\otimes^L_{A_X}L).
\]

It is associative, unital, graded-commutative, and compatible with proper
base change.
\end{definition}

\begin{theorem}[Proper Poincare duality]
If \(X\) is proper as well as smooth, the comparison
\(\RGc(X,K)\cong\RG(X,K)\) identifies \(\beta_K\) with a perfect pairing

\[
 \RG(X,K)\otimes_A^L\RG(X,K^\vee(d))[2d]\longrightarrow A.
\]

Over a field, it gives perfect degreewise pairings

\[
 H^i(X,K)\times H^{2d-i}(X,K^\vee(d))\longrightarrow A.
\]
\end{theorem}

\begin{definition}[Relative trace]
Let \(f:X\to Y\) be smooth and proper of pure relative dimension \(d\).
Smooth purity and the exceptional counit define

\[
 \Tr_f:
 Rf_*A_X(d)[2d]=Rf_!A_X(d)[2d]
 \xrightarrow{\sim}Rf_!Rf^!A_Y\longrightarrow A_Y.
\]
\end{definition}

\begin{definition}[Relative Poincare pairing]
For a perfect \(K\in\Dctf(X,A)\), define

\[
 Rf_*K\otimes^L_{A_Y}Rf_*(K^\vee(d))[2d]
 \longrightarrow A_Y
\]

by the relative cup product, evaluation on \(X\), and \(\Tr_f\).
Explicitly it is

\[
\begin{aligned}
Rf_*K\otimes^LRf_*(K^\vee(d))[2d]
&\xrightarrow{\smile_f}
Rf_*(K\otimes^LK^\vee(d))[2d]\\
&\xrightarrow{Rf_*(\operatorname{ev}_K(d))[2d]}
Rf_*A_X(d)[2d]
\xrightarrow{\Tr_f}A_Y.
\end{aligned}
\]
\end{definition}

\begin{definition}[Perfect pairing over a base scheme]
Let \(P,Q\in D(Y,A)\) be perfect over \(A_Y\).  A morphism

\[
 b:P\otimes^L_{A_Y}Q\longrightarrow A_Y
\]

is perfect if its tensor--Hom adjoint
\(Q\to\RHom_Y(P,A_Y)\) is an isomorphism.
\end{definition}

\begin{theorem}[Relative perfectness of the two direct images]
If \(f:X\to Y\) is smooth and proper and \(K\) is perfect over \(A_X\), then

\[
 Rf_*K,\qquad Rf_*(K^\vee(d))[2d]
\]

are perfect over \(A_Y\).  This assertion is stable under arbitrary base
change on \(Y\) through the proper base-change isomorphism.
\end{theorem}

\begin{theorem}[Relative exceptional tensor--Hom comparison]
Let \(f:X\to Y\) be proper and let \(K\in\Dctf(X,A)\).  Evaluation,
projection, and the counit of \(Rf_!\dashv Rf^!\) define a natural
isomorphism

\[
 \operatorname{RD}_{f,K}:
 Rf_*\RHom_X(K,Rf^!A_Y)
 \xrightarrow{\sim}
 \RHom_Y(Rf_!K,A_Y).
\]

The map is the tensor--Hom adjoint of

\[
 Rf_!K\otimes^LRf_*\RHom_X(K,Rf^!A_Y)
 \longrightarrow Rf_!Rf^!A_Y
 \longrightarrow A_Y.
\]

It is natural in \(K\), compatible with composition of proper morphisms, and
compatible with proper base change.
\end{theorem}

\begin{proposition}[Identification with the relative Poincare adjoint]
If \(f\) is smooth and proper of pure relative dimension \(d\) and \(K\) is
perfect, smooth purity and dualizability give

\[
 \RHom_X(K,Rf^!A_Y)
 \xrightarrow{\sim}K^\vee(d)[2d],
 \qquad Rf_!K\xrightarrow{\sim}Rf_*K.
\]

Under these isomorphisms, \(\operatorname{RD}_{f,K}\) becomes

\[
 Rf_*(K^\vee(d))[2d]
 \longrightarrow\RHom_Y(Rf_*K,A_Y),
\]

and this morphism is exactly the tensor--Hom adjoint of the relative
Poincare pairing defined above.
\end{proposition}

\begin{theorem}[Smooth proper relative duality]
The relative Poincare pairing is perfect.  Its adjoint is the isomorphism

\[
 Rf_*(K^\vee(d))[2d]
 \xrightarrow{\sim}
 \RHom_Y(Rf_*K,A_Y).
\]

For every geometric point \(\bar y\to Y\), proper base change identifies its
fiber with the Poincare pairing on \(X_{\bar y}\).
\end{theorem}

\begin{corollary}[Duality of relative cohomology sheaves]
Assume \(A\) is a field and the cohomology sheaves are lisse.  For every
\(i\), relative duality gives

\[
 \mathcal H^{2d-i}\bigl(Rf_*(K^\vee(d))\bigr)
 \xrightarrow{\sim}
 \mathcal H^i(Rf_*K)^\vee.
\]

This isomorphism commutes with arbitrary base change on \(Y\).
\end{corollary}

\section{Frobenius reciprocity}

Let \(X_0\) be smooth and separated of pure dimension \(d\) over
\(k_0=\mathbf F_q\), put \(\bar k=\overline{\mathbf F}_q\), and let
\(X=X_0\times_{k_0}\bar k\).

\begin{definition}[Galois and scheme Frobenius]
Define arithmetic and geometric Galois Frobenius by

\[
 \sigma:\bar k\longrightarrow\bar k,\quad z\longmapsto z^q,
 \qquad \gamma=\sigma^{-1}.
\]

Base change in the \(\bar k\)-factor gives semilinear inverse automorphisms
\(g_\sigma,g_\gamma:X\to X\).  Let
\(\Phi_0:X_0\to X_0\) be the \(q\)-power scheme Frobenius and put

\[
 \Phi=\Phi_0\times_{\operatorname{Spec}k_0}
 \operatorname{id}_{\operatorname{Spec}\bar k}:X\longrightarrow X.
\]

The map \(\Phi\) is finite, radicial, proper, and a universal homeomorphism;
it need not be an automorphism.  Absolute \(q\)-power Frobenius factors as
\(g_\sigma\Phi=\Phi g_\sigma\).  Topological invariance gives a canonical
natural isomorphism

\[
 \kappa:L\Phi^*\xrightarrow{\sim}Lg_\gamma^*,
\]

and coherent iterates
\(
 \kappa^{(m)}:L(\Phi^m)^*\xrightarrow{\sim}L(g_\gamma^m)^*
\).
\end{definition}

\begin{definition}[Finite-stage strict Weil category]
For a finite coefficient ring \(\Lambda\), let
\(\mathsf{Sh}_{\mathrm W}(X,\Lambda)\) be the abelian category of pairs

\[
 (\mathcal F,\varphi_{\mathcal F}),\qquad
 \varphi_{\mathcal F}:g_\gamma^*\mathcal F\xrightarrow{\sim}\mathcal F,
\]

with equivariant morphisms.  A strict finite-stage Weil complex is a cochain
complex in this abelian category.  Localizing at equivariant
quasi-isomorphisms gives the ordinary triangulated category

\[
 D_{\mathrm W}(X,\Lambda)=
 D\bigl(\mathsf{Sh}_{\mathrm W}(X,\Lambda)\bigr).
\]

The full triangulated subcategory whose underlying complex belongs to
\(\Dctf(X,\Lambda)\) is denoted
\(\DctfW(X,\Lambda)\).  Thus an object includes a specified isomorphism
\[
 \varphi_K:Lg_\gamma^*K\xrightarrow{\sim}K
\]
represented in the equivariant derived category, not merely an isomorphism
of underlying objects.
\end{definition}

\begin{definition}[Adic, rational, and algebraically closed Weil categories]
An object of \(\DctfW(X,\OO_E)\) is a normalized system
\[
 (K_n,\varphi_n,\rho_n)_{n\geq1},
 \qquad (K_n,\varphi_n)\in\DctfW(X,\Lambda_n),
\]
where
\[
 \rho_n:
 K_{n+1}\otimes^L_{\Lambda_{n+1,X}}\Lambda_{n,X}
 \xrightarrow{\sim}K_n
\]
is an equivariant isomorphism: after the canonical exchange of reduction and
\(Lg_\gamma^*\), the square
\[
\begin{CD}
Lg_\gamma^*\red_{n+1,n,X}K_{n+1}
 @>{\red(\varphi_{n+1})}>> \red_{n+1,n,X}K_{n+1}\\
@V{Lg_\gamma^*\rho_n}V{\sim}V @VV{\rho_n}V\\
Lg_\gamma^*K_n @>{\varphi_n}>> K_n
\end{CD}
\]
commutes.  The underlying system has common cohomological bounds, a common
constructible stratification, and one common finite Tor-amplitude interval.
Morphisms are compatible systems of equivariant derived morphisms.

Define the rational and algebraically closed coefficient categories by
\[
 \DctfW(X,E)=
 \operatorname{Kar}\bigl(\DctfW(X,\OO_E)[\varpi^{-1}]\bigr),
\]
\[
 \DctfW(X,\Qlb)=
 \operatorname{Kar}\!\left(
  2\!\operatorname*{colim}_{E/\Ql\ \mathrm{finite}}\DctfW(X,E)
 \right).
\]
An object at the \(E\)-stage is represented by a rationalized adic
Weil system and an idempotent; an object at the \(\Qlb\)-stage is represented
by such data over some finite \(E/\Ql\), modulo further finite extension.
There are exact forgetful functors
\[
 \operatorname{For}_A:\DctfW(X,A)\longrightarrow\Dctf(X,A)
\]
at all four coefficient stages, compatible with reduction, scalar extension,
and idempotent splittings.  A Weil complex below means an object
\((K,\varphi_K)\in\DctfW(X,A)\), and \(K\) denotes its underlying complex.
\end{definition}

\begin{definition}[Iterated Weil and scheme-Frobenius maps]
For \(m\geq1\), define

\[
 \varphi_K^{[m]}:
 L(g_\gamma^m)^*K\xrightarrow{\sim}K,
 \qquad
 \varphi_K^{[m]}=
 \varphi_K\circ Lg_\gamma^*\varphi_K\circ\cdots\circ
 L(g_\gamma^{m-1})^*\varphi_K.
\]

Its scheme-Frobenius form is

\[
 u_m=\varphi_K^{[m]}\circ\kappa_K^{(m)}:
 L(\Phi^m)^*K\longrightarrow K.
\]

The iterates satisfy
\(
 \varphi_K^{[m+n]}=
 \varphi_K^{[m]}\circ L(g_\gamma^m)^*\varphi_K^{[n]}
\)
under the canonical pullback identifications.
\end{definition}

\begin{definition}[Geometric-Frobenius operators on global sections]
Base change for the automorphism \(g_\gamma^m\) gives canonical maps

\[
 g_{\gamma}^{m,*}:\RG(X,K)\xrightarrow{\sim}
 \RG\bigl(X,L(g_\gamma^m)^*K\bigr),
\]

\[
 g_{\gamma,c}^{m,*}:\RGc(X,K)\xrightarrow{\sim}
 \RGc\bigl(X,L(g_\gamma^m)^*K\bigr).
\]

Define

\[
 F_K^{(m)}=\RG(X,\varphi_K^{[m]})\circ g_\gamma^{m,*},
 \qquad
 F_{c,K}^{(m)}=\RGc(X,\varphi_K^{[m]})\circ g_{\gamma,c}^{m,*}.
\]

These are typed endomorphisms of \(\RG(X,K)\) and \(\RGc(X,K)\),
respectively.  Put
\[
 F_K=F_K^{(1)},\qquad F_{c,K}=F_{c,K}^{(1)},
\]
and write
\[
 F_{K,i}=H^i(F_K),\qquad F_{c,K,i}=H^i(F_{c,K})
\]
for the induced endomorphisms of \(H^i(X,K)\) and \(H_c^i(X,K)\).
When the cohomological degree is fixed, \(F_K\) and \(F_{c,K}\) also denote
these induced endomorphisms; \(F_{K^\vee(r)}\) denotes the operator obtained
from the dual-and-Tate datum defined below.
Via \(\kappa^{(m)}\), the same operators are obtained from the
scheme-Frobenius maps \(u_m\).  Reduction, rationalization, finite extension
of coefficients, and passage to an idempotent summand carry these operators
to the corresponding operators at the next coefficient stage.
\end{definition}

\begin{definition}[Dual and Tate Weil data]
Let
\[
 \varphi_A:Lg_\gamma^*A_X\xrightarrow{\sim}A_X
\]
be the canonical Weil structure on the constant coefficient object.  For a
perfect Weil complex \((K,\varphi_K)\), define the contragredient datum on
\(K^\vee=\RHom_X(K,A_X)\) by the typed composite
\[
\begin{split}
 \varphi_{K^\vee}:Lg_\gamma^*\RHom_X(K,A_X)
 &\xrightarrow{\sim}
 \RHom_X(Lg_\gamma^*K,Lg_\gamma^*A_X)\\
 &\xrightarrow{\RHom(\operatorname{id},\varphi_A)}
 \RHom_X(Lg_\gamma^*K,A_X)\\
 &\xrightarrow{\RHom(\varphi_K^{-1},\operatorname{id})}
 \RHom_X(K,A_X).
\end{split}
\]

Let \(\tau_1:L g_\gamma^*A_X(1)\xrightarrow{\sim}A_X(1)\) be the canonical
Tate datum with geometric-Frobenius scalar \(q^{-1}\).  Tensor powers and
tensor inverses define
\[
 \tau_r:L g_\gamma^*A_X(r)\xrightarrow{\sim}A_X(r)
 \qquad(r\in\mathbf Z).
\]
The tensor product of \(\varphi_{K^\vee}\) and \(\tau_r\) is denoted
\(\varphi_{K^\vee(r)}\).  In particular, geometric Frobenius acts on
\(A_X(d)\) by \(q^{-d}\) and on \(A_X(-d)\) by \(q^d\).
\end{definition}

\begin{proposition}[Dual and Tate data at every coefficient stage]
The datum \(\varphi_{K^\vee}\) defines an object
\((K^\vee,\varphi_{K^\vee})\in\DctfW(X,A)\), coefficient evaluation
\[
 K\otimes^LK^\vee\longrightarrow A_X
\]
is a morphism of Weil objects, and every
\((K^\vee(r),\varphi_{K^\vee(r)})\) belongs to \(\DctfW(X,A)\).  These
objects and morphisms commute with finite reduction, adic normalization,
rationalization, finite coefficient extension, and idempotent splitting.
\end{proposition}

\begin{theorem}[Weil compatibility of the ingredients]
Mixed-support cup product, coefficient evaluation, smooth purity, and the
normalized trace commute with the specified Weil isomorphisms.  The derived
Poincare pairing is a morphism in the corresponding coefficient-stage Weil
category.
\end{theorem}

\begin{theorem}[Equivariance of the twisted pairing]
Let \((K,\varphi_K)\) be a perfect Weil complex.  With
\(F_{K^\vee(d)}\) formed from the datum \(\varphi_{K^\vee(d)}\), the
cohomological operators satisfy

\[
 \bigl\langle F_{c,K}x,F_{K^\vee(d)}y\bigr\rangle_i
 =\langle x,y\rangle_i
\]

for
\(x\in H_c^i(X,K)\) and
\(y\in H^{2d-i}(X,K^\vee(d))\), where the action on the second factor
includes the scalar \(q^{-d}\) on \(A(d)\).
\end{theorem}

\begin{corollary}[Untwisted similitude]
The untwisted pairing takes values in \(A(-d)\), on which geometric
Frobenius acts by \(q^d\), and it satisfies

\[
 \bigl\langle F_{c,K}x,F_{K^\vee}y\bigr\rangle_i^{\mathrm{untw}}
 =q^d\langle x,y\rangle_i^{\mathrm{untw}}.
\]
\end{corollary}

\begin{corollary}[Reciprocal eigenvalues]
For field coefficients, the operator \(F_{K^\vee}\) on
\(H^{2d-i}(X,K^\vee)\) is conjugate to

\[
 q^d(F_{c,K}^{-1})^{\mathsf t}
\]

relative to the perfect pairing.  Thus an eigenvalue \(\alpha\) on
\(H_c^i(X,K)\) corresponds, with the same algebraic multiplicity, to
\(q^d/\alpha\) on \(H^{2d-i}(X,K^\vee)\).  If \(X\) is proper, the first
space is \(H^i(X,K)\).
\end{corollary}

\section{Coefficient-change rectangles}

\begin{theorem}[Finite reduction]
For \(m\leq n\), finite reduction commutes with ordinary and compactly
supported global sections, coefficient duals, Tate twists, cup products,
smooth purity, evaluation, and trace.  In particular, the square

\[
\begin{CD}
\red\RG(X,K^\vee(d))[2d] @>{\red(\beta_K^\sharp)}>>
\red R\operatorname{Hom}(\RGc(X,K),\Lambda_n)\\
@V{\sim}VV @VV{\sim}V\\
\RG(X,(\red K)^\vee(d))[2d] @>{\beta_{\red K}^\sharp}>>
R\operatorname{Hom}(\RGc(X,\red K),\Lambda_m)
\end{CD}
\]

commutes, where \(\beta^\sharp\) denotes the adjoint of the pairing.
\end{theorem}

\begin{theorem}[Adic reconstruction and scalar extension]
The finite-level pairings of a normalized adic object form a morphism of
normalized systems and induce the \(\OO_E\)-valued Poincare pairing.
Rationalization and finite coefficient-field extension carry this pairing,
its adjoint, every degree-indexing map, the relative pairing, and all
Frobenius squares to their counterparts over the extended coefficient field.
\end{theorem}

\begin{theorem}[\(\Qlb\)-coefficient pairing]
Represent \(K\) and its idempotent over a finite coefficient field.  The
Poincare pairing preserves the corresponding direct summands and is
independent of every further finite extension.  It defines the
perfect derived and degreewise pairings over \(\Qlb\).
\end{theorem}

\end{document}
